\chapter{NETRA}
\label{ch:u-netra}

NETRA is the command center: open it in a browser on the laptop that runs it
(\url{http://localhost}) or on any device on the same network, and log in.

\section{What is where}

\begin{center}
\begin{tabularx}{\linewidth}{@{}>{\bfseries}p{0.3\linewidth}X@{}}
  \toprule
  Page & For \\
  \midrule
  Dashboard & every robot at a glance; what needs attention (e-stop, alarms, battery, lost);
    \ui{E-stop} and its release \\
  Operations map & where the robots are; a robot's card: battery, mode, health, REC, and
    \menupath{Live control} (the recorder, the robot's computer) \\
  Patrols & every patrol with its recap and timeline; \ui{Start a patrol now} \\
  Schedules & patrols by day and time \\
  Routes & the routes the robot taught \\
  Substation & overview, findings, equipment, inspection points \\
  Reports & PDF and Excel exports \\
  \bottomrule
\end{tabularx}
\end{center}

\section{Alerts}

The robot reports when it is blocked, off its route or failed, and when its e-stop engages.
NETRA shows them on the dashboard and, where configured, sends them to Telegram. Every finding
of warning or critical severity appears under \menupath{Substation › Findings}.

\section{Commands to the robot}

From the robot's card: pause, resume, stop the patrol, localize, the e-stop and its release, and
the recorder. The robot answers each one; a refused command shows the robot's reason.
